\section{Giriş}\label{sec:intro}

Elma hasadı, yoğun iş gücü gerektiren seçmeli bir işlemdir: aynı ağaçtaki meyveler eşzamanlı olgunlaşmaz,
çeşitler farklı tarihlerde olgunluğa erişir ve taze tüketime yönelik bahçelerde toplama çoğunlukla
iki ya da daha fazla geçişte yapılır \citep{simnitt_2023_advancements,hall_2001_effect}. Robotik
hasat makineleri günümüzde ticari bahçelerde elma toplamaktadır \citep{silwal_2017_design,
zhang_2024_automated}; yer platformlarına kabloyla bağlanan drone'ları da içeren toplama robotu
filoları \citep{oitzman_2021_tevel} ise geliştirilme aşamasındadır. Böyle bir filonun ağaç düzeyinde
koordinasyonu bir rotalama problemidir: olgun meyvesi bulunan her ağaç bir robota
atanır, her robot baş tarladaki bir istasyondan ayrılır ve oraya geri döner;
toplama, son robot geri döndüğünde tamamlanır. Bu tamamlanma süresinin (makespan) en aza indirilmesi,
min--maks çoklu gezgin satıcı problemine (min--maks mTSP) ve bu problemin
kapasiteli ve çok seferli varyantlarına karşılık gelir \citep{bektas2006,bochtis2009}. Birkaç
hektarlık bir bahçede binlerce ağaç bulunur ve plan uzun süre geçerli kalmaz: ağaç başına
olgun meyve miktarı toplamadan önce yalnızca kabaca bilinir, robotlar çalıştıkça kameralar
kestirimleri iyileştirir ve plan, bir robot her döndüğünde yeniden hesaplanır.
Her yeniden hesaplamanın, robot yeniden yola çıkmaya hazır olmadan önce tamamlanması gerekir.

Kullanılabilecek iki çözücü ailesi bulunmaktadır. Min--maks mTSP için tasarlanmış
metasezgiseller \citep{hehao2022,hehao2023,hehaoxia2024,zheng2022,mahmoudinazlou2024} kıyaslama
kümelerinde bilinen en iyi çözümlere ulaşmaktadır; genel amaçlı çözücü
LKH-3 \citep{helsgaun2017} ise uygulamada olağan referans olarak kullanılmaktadır. Her iki aile de
örnek başına dakikalardan saatlere uzanan bütçelerle değerlendirilmekte ve nihai değerleriyle raporlanmaktadır;
bildiğimiz kadarıyla literatürde bu yöntemlerin ilk saniyelerdeki, bahçe büyüklüğüne ve yapısına sahip örneklerdeki
davranışı ölçülmemiştir. Bahçeler probleme geometrik nitelikte ikinci bir boyut ekler:
robotlar, yönelimi önceden bilinmeyen ve kafes düzeninde düzensizlikler barındıran bir dikimin
sıra aralarında hareket eder.

Bu makalenin katkıları şunlardır:
\begin{enumerate}
\item \textbf{PSB} (Path--Split--Balance): depolu Öklidyen min--maks mTSP için
önce rota, sonra küme yaklaşımına dayalı bir işlem hattı: Greedy-Edge dev turu, optimal
min--maks bölme (Lemma~\ref{lem:split}), cilalama, dengeleme ve her an sonuç veren
(anytime) geniş komşuluk araması (LNS); ayrıca yöntemin kasa kapasiteli seferler, ağaç konumlarından
kestirilen bahçe ekseninde oluşturulan dev tur ve olay tetiklemeli yeniden planlama ile
bahçelere genişletilmesi.
\item \InstAll{} örnek (mTSPLib, sentetik bahçeler, OpenStreetMap'ten alınan on altı gerçek
bahçe sınırı, VLSI) ve büyük literatür kümelerindeki \LitCount{} örnek üzerinde, üç
yapılandırmadaki LKH-3 ile herkese açık uygulaması bulunan, literatürde öne çıkan dört
min--maks mTSP koduna (İlgili çalışmalar bölümünde tanıtılan HSNR, MA, MILS, HGA) karşı yürütülen, ablasyon ve tekrarlı koşular içeren,
zaman boyutunda çözümlenmiş, tek çekirdekli bir kıyaslama. Bu kıyaslama, hangi yöntemin hangi durumda
tercih edilmesi gerektiğini ortaya koymaktadır: ilk saniyelerden dakikalara uzanan aralıkta ve
binlerce hedef içeren örneklerde PSB, orta büyüklükteki örnekler için birkaç dakikalık süre
bulunduğunda ise MILS.
\item Kablolu drone platformlarıyla taze tüketime yönelik seçmeli elma hasadının, parametreleri
tarım literatüründen alınan ve gerçek sınırlar üzerinde kusurlu geometriye sahip
kare dikimli ve Tall Spindle bahçelerde yürütülen ayrık olaylı benzetimi.
Bu benzetim, PSB ile yeniden planlamayı sabit sıra bantlarına dayalı basit bir statik referansla,
açgözlü dağıtımla, OR-Tools ile ve uç devralmalı merkezi olmayan bir $k$-d bölüntüsüyle karşılaştırmakta,
bahçe ekseni kestiriminin katkısını ve planlama gecikmesinin etkisini
nicel olarak ortaya koymakta ve bir duyarlılık analizi içermektedir.
\item Her sayıyı, tabloyu ve şekli üreten tüm kodlar, kayıtlar ve betikler ile
benzetimi yapılan görevlerin Blender ortamında hazırlanmış yeniden oynatımı.
\end{enumerate}

Makalede önce ilgili çalışmalar gözden geçirilmiş ve problem tanımlanmış, ardından PSB ile
genişletmeleri, kıyaslama düzeneği ve sonuçları ve hasat benzetimi açıklanmış;
son olarak çıkarımlar ve sınırlılıklar tartışılmıştır.
